\scanpage{341}
% @ top centre, overwritten numeral
\mspage{6\deleted{\unsure{4}}5}

% @ top paragraph, first two lines
je ne me suis pas suffisamment attaché à développer des exemples et des \unsure{contr'exemples|contre'exemples}.

% @ upper third, observations and hypothetical implication
Sans compter qu'il me semble que tous les loc. fond. \inserted{sur $\mathrm{Cat}$} qu'on rencontre ``en pratique'' satisfont déjà à $\mathrm{Loc}(8)$ et a fortiori à $\mathrm{Loc}(8\mathrm{bis})$, donc pour eux on a \deleted{\unsure{bien}} l'implication escomptée. La \unsure{$=$|même} remarque s'applique à l'implication hypothétique

% @ centre, question mark over arrow
\[
 D_{W_0}\overset{?}{\Longrightarrow}A,
\]
% @ centre, qualification
sans hypothèse $\mathrm{Loc}(8\ \mathrm{bis})$ sur $W_0$ et sur $W$.

% @ lower middle through foot, candidate and next question
\deleted{Il faudrait \unsure{exp}} Un candidat évident pour un contre-exemple déjà à la première question serait $W_\infty={}$ plus petit localiseur fondamental sur $\mathrm{Cat}$, duquel il est presque évident qu'il ne satisfait pas à $\mathrm{Loc}(8\deleted{\unsure{\mathrm{bis}}})$ \unsure{ni} même $\mathrm{Loc}(8\mathrm{bis})$. Mais il n'y a aucune raison qu'il satisfasse \inserted{non plus} à $F$, \deleted{donc} de telle façon que (s'il n'y satisfait pas\unsure{\deleted{)}|)}, donc non plus à $D$) l'implication \inserted{$D\Rightarrow A$} est bel et bien valable. Cela nous amène à la question, plus substantielle, si on n'aurait pas \deleted{\unsure{cela}} l'implication
% archive sheet number: 340
